% Generated draft: requires mathematical and source review.
\documentclass[11pt]{article}
\usepackage[margin=25mm]{geometry}
\usepackage{fontspec}
\setmainfont{Latin Modern Roman}
\usepackage{amsmath,amssymb,mathtools,unicode-math}
\setmathfont{Latin Modern Math}
\usepackage{enumitem,needspace,xurl,hyperref}
\setlength{\parindent}{0pt}
\setlength{\parskip}{5pt}
\setlength{\emergencystretch}{3em}
\newcommand{\N}{\mathbb{N}}
\newcommand{\Z}{\mathbb{Z}}
\newcommand{\Q}{\mathbb{Q}}
\newcommand{\R}{\mathbb{R}}
\newcommand{\C}{\mathbb{C}}
\newcommand{\F}{\mathbb{F}}
\newcommand{\A}{\mathbb{A}}
\newcommand{\PP}{\mathbb{P}}
\newcommand{\E}{\mathbb{E}}
\newcommand{\eps}{\varepsilon}
\newcommand{\dd}{\,\mathrm{d}}
\begin{document}
\section*{UnsolvedMath ID: 30002688}
\textbf{Source problem code:} catalogue entry
\subsection{Definitions and mathematical statement}
Put \(\mathfrak n=\operatorname{Lie}(N)\), and write
\[
\rho(u)=\operatorname{Ad}(u)|_{\mathfrak n},\qquad H=\rho(U).
\]
This notation presupposes an ambient Lie group in which the adjoint action is defined and satisfies \(\operatorname{Ad}(u)\mathfrak n=\mathfrak n\) for every \(u\in U\). The supplied record does not specify that ambient group or the hypotheses on \(U\) and \(N\).

To interpret transpose in the usual real Riemannian setting, one must also specify a positive-definite inner product \(\langle\cdot,\cdot\rangle\) on the finite-dimensional real vector space \(\mathfrak n\). For such a specified inner product, the transpose \(A^{\mathsf T}\) of an endomorphism \(A\) is characterized by
\[
\langle Ax,y\rangle=\langle x,A^{\mathsf T}y\rangle
\qquad(x,y\in\mathfrak n).
\]
The choice of real setting and inner product here explains the usual meaning of the question; neither is explicitly specified in the supplied statement.

\textbf{Question, with unresolved source hypotheses.} For the groups and inner product intended by the source, is the restricted adjoint image invariant under transpose, that is,
\[
H^{\mathsf T}:=\{A^{\mathsf T}:A\in H\}=H?
\]
Equivalently, is it true that
\[
\forall u\in U\;\exists v\in U:\qquad
\bigl(\operatorname{Ad}(u)|_{\mathfrak n}\bigr)^{\mathsf T}
=\operatorname{Ad}(v)|_{\mathfrak n}?
\]
Here ``self-adjoint action'' means invariance of the image under transpose. It does not require each image element to equal its own transpose, and ``closed under transpose'' does not assert topological closedness.

The supplied background describes \(\mathfrak n\) as a nilradical and places this question in work toward the strong Alekseevskii conjecture. It does not identify the Lie algebra whose nilradical is intended, the relevant geometric assumptions, or the metric determining transpose. Thus the displayed condition gives the precise property being asked about, but the intended class of objects over which it is to hold cannot be recovered from this record. In particular, the record does not determine whether the inner product is fixed by the geometry or is to be chosen existentially. The supplied status is \emph{open}; this draft does not independently verify that status.
\subsection{Short English statement}
Does transposing any transformation in the restricted adjoint image of \(U\) on \(\operatorname{Lie}(N)\) produce another transformation in the same image? The ambient groups, geometric hypotheses, and inner product defining transpose are missing, so the intended scope remains unresolved. The supplied background presents this as a structural question toward the strong Alekseevskii conjecture.
\subsection{Sources}
\begin{itemize}
\item Supplied local UnsolvedMath registry record, permanent ID \(30002688\), source problem code \texttt{catalogue entry}, collection \emph{Oberwolfach Reports Open Problems}. No individual proposer is named; the supplied proposed year is 2014.
\item Supplied source citation: \emph{Mini-Workshop: Einstein Metrics, Ricci Solitons and Ricci Flow under Symmetry Assumptions} (2015), \emph{Oberwolfach Reports}, pp.\ 2529--2568, DOI \texttt{10.4171/owr/2014/45}. \url{https://doi.org/10.4171/owr/2014/45}. The supplied literature list labels this source \emph{Oberwolfach Report 45/2014}. The citation's 2015 date and the record's 2014 report/proposal dating are retained without reconciliation.
\item Supplied literature reference: Bohm--Lafuente, \emph{Non-compact Einstein Manifolds with Symmetry}. \url{https://doi.org/10.1090/jams/1022}.
\item Literature assessment checked at \texttt{2026-08-21}: supplied dated curation, not independently verified. It labels the problem open, describes self-adjointness of the \(U\)-action on the nilradical as a structural step toward the strong Alekseevskii conjecture, and reports that no general theorem resolving this exact action property was found. The supplied background calls this a dated literature triage derived from OWR curation and requests independent verification of time-sensitive status and cited claims.
\item The registry labels the statement \texttt{corrected\_verified} and records a manual semantic, typographic, and LaTeX review claiming that the cleanup is self-contained and adds no unsupported source-dependent content. These are supplied editorial assertions; they do not resolve the missing mathematical hypotheses identified here. No cited source was independently consulted for this draft.
\end{itemize}
\end{document}
